\documentclass[journal]{IEEEtran}

\usepackage{graphicx}
\usepackage{booktabs}
\usepackage{array}
\usepackage{amsmath,amssymb}
\usepackage{url}
\usepackage{microtype}
\usepackage[hidelinks]{hyperref}

% WellPulse IEEE Revival -- R7b source skeleton.
% NEW SOURCE IDENTITY. Does not overwrite rejected TNSM release.
% Governing science: experiments/WP-RT01/R6_INTEGRATED_SCIENTIFIC_ADJUDICATION.md

\begin{document}

\title{Beyond Reconnection: Measuring Recovery Observables and Replay-Induced Continuity Gaps in MQTT Telemetry}

\author{Ahmed Ayoub%
\thanks{Author affiliation and correspondence metadata are intentionally deferred to final editorial reconciliation in R7e; R7b does not alter submission metadata.}}

\maketitle

\begin{abstract}
Telemetry recovery is often summarized by a single reconnect or catch-up time, although network restoration, MQTT session restoration, broker-side QoS-1 acknowledgement closure, receiver-side historical completion, live-data continuity, and stale-record carryover are different observables. This study defines an experimental measurement protocol that separates these states and applies it to MQTT telemetry on FIT IoT-LAB and POWDER. In the tested W1 implementation, serialized backlog replay preserved eventual 10,000/10,000 receiver completeness but repeatedly suspended new record generation. Six historical FIT runs showed sequence-5000-to-5001 gaps of 68.785--70.264~s (mean 69.217~s); three prospective scored runs showed 74.798--76.308~s (mean 75.438~s). The prospective runs directly instrumented sender PUBACK closure and receiver completion, but 0/3 produced either a causal witness or a clock-resolved positive M2-to-M3 separation, so that interval remains unresolved. A POWDER trace independently exhibited stale carryover: newer records arrived before an older record after restoration. The results support recovery-observable measurement and identify an implementation-specific completeness--continuity trade-off without claiming a new persistence architecture, a universal MQTT effect, or pooled cross-testbed inference.
\end{abstract}

\begin{IEEEkeywords}
MQTT, telemetry recovery, failure recovery, data durability, backlog replay, continuity, receiver reconciliation, fault injection, IoT testbeds.
\end{IEEEkeywords}

\section{Introduction}
\label{sec:introduction}

% R7d: rebuild motivation around the overloaded term recovery.
% Do not lead with architecture novelty or repeated M2-to-M3 misclassification.

\subsection{Research Questions}
This revival is organized around four bounded research questions:
\begin{enumerate}
    \item[\textbf{RQ1}] Which experimentally observable endpoints are required to distinguish network-path restoration, MQTT-session restoration, broker-side QoS-1 acknowledgement closure, receiver-side historical completion, live-data continuity, and receiver-side stale carryover?
    \item[\textbf{RQ2}] Under the tested serialized W1 backlog-replay implementation, what live-generation discontinuity is observed in the historical and prospective FIT IoT-LAB campaigns?
    \item[\textbf{RQ3}] When M2 and M3 are directly instrumented prospectively, does the qualified evidence resolve a positive M2-to-M3 separation under the frozen causal-witness and clock-uncertainty rules?
    \item[\textbf{RQ4}] What complementary, bounded evidence does POWDER provide about failure-domain-specific recovery and receiver-side stale carryover without pooling the two testbeds?
\end{enumerate}

\subsection{Contributions}
The reconstructed paper is limited to four contributions:
\begin{enumerate}
    \item[\textbf{A1}] \textbf{Recovery-observable measurement protocol.} A reproducible measurement discipline that separates M0--M5 and prevents reconnect, session restoration, PUBACK closure, receiver completion, continuity, and ordering from being treated as interchangeable evidence.
    \item[\textbf{A2}] \textbf{Repeated W1 completeness--continuity result.} Across separate historical and prospective FIT campaigns, serialized backlog replay repeatedly produced a tens-of-seconds suspension of new record generation while eventual receiver completeness was preserved.
    \item[\textbf{A3}] \textbf{Cross-environment triangulation without pooling.} FIT provides implementation-specific replay evidence; POWDER provides complementary bounded physical failure-domain and stale-carryover evidence. No cross-platform effect estimate is constructed.
    \item[\textbf{A4}] \textbf{Explicit negative/bounded result.} The prospective Stage-A campaign produced 0/3 positive M2-to-M3 results and 3/3 UNRESOLVED outcomes; the negative result is retained rather than converted into a favorable recovery claim.
\end{enumerate}

% R7d: add final Introduction roadmap. Single-author voice only.

\section{Related Work and Claim Boundary}
\label{sec:related}

% R7d block 1: MQTT persistence/session/store-and-forward.
% R7d block 2: robustness, broker performance, fault injection.
% R7d block 3: layered network/service/data continuity and subscriber confirmation.
% R7d block 4: retained bounded measurement gap.
% Prohibited: first, to our knowledge, novel recovery architecture.

\section{Recovery Observables and Experimental Protocol}
\label{sec:protocol}

\subsection{Recovery-observable taxonomy}

\begin{table*}[t]
\centering
\caption{Frozen recovery observables used by the revival study.}
\label{tab:observables}
\begin{tabular}{p{0.06\textwidth}p{0.22\textwidth}p{0.31\textwidth}p{0.31\textwidth}}
\toprule
\textbf{ID} & \textbf{Observable} & \textbf{Operational event / authority} & \textbf{Does not establish} \\
\midrule
M0 & Network-path restoration & First successful post-restoration user-plane probe under the frozen probing rule & MQTT session restoration, broker acceptance, or receiver delivery \\
M1 & MQTT-session restoration & Publisher-side successful CONNACK & Historical completion or subscriber receipt \\
M2 & Sender-to-broker QoS-1 closure & Publisher timestamp of the final PUBACK covering the frozen historical set $H$ & Subscriber receipt or receiver-side historical completion \\
M3 & Receiver historical-state completion & Receiver timestamp at which every unique record in $H$ has been observed & Live-generation continuity or freshness \\
M4 & Live-data continuity & Source-generation timing, including the replay-associated generation gap & Receiver completion latency \\
M5 & Receiver ordering / stale carryover & Same-receiver arrival order and older-after-newer evidence & Prevalence or probability across deployments \\
\bottomrule
\end{tabular}
\end{table*}

\subsection{Historical FIT campaign}
% R7c:
% B0 is a non-durable mechanism-isolation control only.
% Historical W1 generation gaps:
% 68.992, 69.466, 68.944, 68.785, 70.264, 68.852 s.
% Mean 69.217 s; range 68.785--70.264 s.
% Historical 68--70 s must not be called receiver M3 completion.

\subsection{Prospective FIT Stage-A campaign}
% H = sequences 3001--5000, exactly 2000 records.
% Same-receiver causal witness: seq5001 before any remaining H record.
% Conservative pre/post cross-host clock-offset interval.
% Actual outcome: 0/3 positive; Stage B not executed.

\subsection{POWDER physical campaign}
% Preserve named failure domains and exact/censored/upper-bound semantics.
% Add E3 cycle-3 stale-carryover existence observation:
% seq251--255 arrived before older seq231 on the same receiver clock.
% Do not infer prevalence.

\subsection{Evidence separation and non-pooling}
% Historical FIT, prospective FIT, and POWDER remain separate evidence classes.

\section{Results}
\label{sec:results}

\subsection{Historical FIT: eventual completeness and replay-associated continuity gap}
% R7c full prose. Headline is M4 / implementation behavior.

\subsection{Prospective FIT: direct M2/M3 instrumentation and Stage-A adjudication}

\begin{table*}[t]
\centering
\caption{Prospective scored FIT Stage-A results. All three runs were valid.}
\label{tab:stagea}
\begin{tabular}{lrrrrcl}
\toprule
\textbf{Run} &
\textbf{Seq. 5000$\rightarrow$5001 gap (s)} &
\textbf{$H$ PUBACK} &
\textbf{$H$ receiver} &
\textbf{Causal witness} &
\textbf{Conservative $D_{23}$ interval (s)} &
\textbf{Class} \\
\midrule
R1 & 75.208 & 2000/2000 & 2000/2000 & No & $[-5.038,\,+4.801]$ & UNRESOLVED \\
R2 & 76.308 & 2000/2000 & 2000/2000 & No & $[-4.857,\,+4.932]$ & UNRESOLVED \\
R3 & 74.798 & 2000/2000 & 2000/2000 & No & $[-4.858,\,+4.815]$ & UNRESOLVED \\
\bottomrule
\end{tabular}
\end{table*}

% State 0/3 positive; decision NOT_CONFIRMED_STOP.
% Raw cross-clock differences are not ordering facts because intervals cross zero.

\subsection{POWDER: bounded failure-domain and stale-carryover evidence}
% R7c full prose.

\subsection{Cross-evidence synthesis}
% Common result is measurement discipline, not a pooled effect.

\section{Discussion}
\label{sec:discussion}
% R7d full prose.

\section{Threats to Validity}
\label{sec:threats}
% R7d: historical endpoint error, prospective clock uncertainty, no causal
% witness, 0/3 does not prove equality, W1 specificity, B0 asymmetry,
% three run-level replicates, non-pooling, bounded POWDER observation.

\section{Reproducibility and Data Availability}
\label{sec:repro}
% R5 raw authority SHA-256:
% 5b52005e48d8987091c447186eb39261824f32df15653c3d2773a031b062d4a2
% Public artifact must remain sanitized.

\section{Conclusion}
\label{sec:conclusion}
% R7d full rewrite under R6 claim ceiling.

% R7e will rebuild the complete IEEE-formatted reference list.
\begin{thebibliography}{1}
\bibitem{mqtt5}
A.~Banks, E.~Briggs, K.~Borgendale, and R.~Gupta, Eds.,
MQTT Version 5.0, OASIS Standard, Mar. 2019.
\end{thebibliography}

\end{document}
